\begin{Thanks}
    \text 歷經兩年的研究所學習與論文撰寫，本論文終於得以完成。回首這段旅程，自香港來臺求學至今已歷經數年，從大學就讀資訊管理學系，到升讀國立臺北科技大學資訊與財金管理系碩士班，再到完成碩士學位，這一路上有許多師長、同事、家人及朋友的鼓勵與支持，才能讓我順利完成人生的重要階段。謹藉此機會，向所有曾經幫助過我的師長、同事及親友，獻上最誠摯的感謝。
    
    \text 首先，我要衷心感謝我的指導教授 彭祖乙教授。老師於研究期間，不僅在研究方向、論文架構及學術方法上悉心指導，更在我遭遇研究瓶頸時耐心引導，給予許多寶貴建議，使我逐步建立研究思維，培養獨立分析與解決問題的能力。老師嚴謹的治學態度、對研究品質的要求，以及鼓勵學生勇於探索與持續精進的精神，讓我受益匪淺，也將成為我未來無論在職涯或研究道路上的重要養分。

    \text 同時，誠摯感謝本次論文口試委員 林敬皇教授 與 胡念祖教授 百忙之中撥冗審閱本論文並擔任口試委員。特別感謝 林敬皇教授，除了於本次論文口試提供許多寶貴建議外，更於論文計畫書提案口試階段即給予我許多研究方向與內容上的指導，讓我重新檢視研究架構，並持續修正研究方向，使本論文得以更加完整與扎實。亦感謝來自國立虎尾科技大學的 胡念祖教授，於口試過程中提出許多具建設性的建議與不同觀點，讓我得以從更全面的角度重新思考研究內容，也使本論文在完整性及未來研究方向上獲得更多啟發。

    \text 此外，我要特別感謝 資誠創新諮詢有限公司（PwC Taiwan—Consulting / Cloud Innovation Center）。自大學期間有幸加入公司擔任實習生以來，公司便提供我寶貴的實務學習機會；大學畢業後升讀研究所期間，亦讓我持續以約聘人員的身分參與團隊工作，使我能夠兼顧學業與工作，在學術研究與實務經驗之間取得良好的平衡。衷心感謝管理顧問及數位轉型核心團隊所有主管與同仁，在研究所兩年間對我的支持、信任、包容與鼓勵，使我得以持續累積顧問實務經驗，也讓本研究能夠結合更多產業實務的思維。如今順利完成碩士學業，亦有幸以正式員工的身分繼續服務於公司，感謝公司一路以來的信任與栽培，給予我持續學習與發展的機會，也讓我對未來的職涯充滿期待。

    \text 最後，我也要特別感謝遠在香港的家人、朋友，以及一路上所有關心與支持我的人。感謝你們始終尊重我的每一個選擇，支持我離鄉來臺求學，在我面對課業、研究與工作的壓力時，始終給予我最大的鼓勵與信任。正因為有你們的陪伴與鼓勵，讓我在兼顧研究所課業、論文研究與工作之餘，始終能夠堅持初心，勇敢面對每一次挑戰，也讓我順利完成這段充實而難忘的求學旅程。這份成果，同樣屬於你們。

    \text 論文的完成並非終點，而是另一段學習與成長旅程的開始。未來，我將持續秉持在研究所期間所培養的研究精神與專業態度，不斷精進自我，將所學應用於資訊科技、人工智慧、數位轉型與企業管理等實務領域，並以更謙遜的態度持續學習，迎接未來更多挑戰。謹以此文，向所有曾經陪伴我、支持我、鼓勵我的每一位，致上最誠摯的感謝。
\end{Thanks}